% Supplementary Material — compiles standalone with the AAAI-27 author kit
% (place next to aaai2027.sty). Anonymous: contains no author-identifying info.
\documentclass[letterpaper]{article}
\usepackage[submission]{aaai2027}
\usepackage[hyphens]{url}
\usepackage{graphicx}
\urlstyle{rm}
\def\UrlFont{\rm}
\usepackage{natbib}
\usepackage{caption}
\frenchspacing
\usepackage{booktabs}
\usepackage{amssymb}
\setcounter{secnumdepth}{1}
\renewcommand\thesection{S\arabic{section}}
\renewcommand\thetable{S\arabic{table}}
\pdfinfo{/TemplateVersion (2027.1)}

\title{Supplementary Material for\\ ``From Prediction to Policy: A Learning-Based Microsimulation of Employee Attrition and the Wage-Elasticity of Retention''}
\author{Anonymous submission}
\affiliations{~}

\begin{document}
\maketitle

\section{Uncertainty Analysis}

\paragraph{Simulation dispersion}
The microsimulation is effectively deterministic at cohort scale. Repeating the
backtest of the main paper (sub-models trained on the 2017--2020 waves,
simulation initialized with the 2021 employee cohort) over ten seeds, the
annual separation rate varies with a standard deviation of at most 0.002
(Table~\ref{tab:sim_dispersion}), and end-year (2024) mean real income varies
by 0.36 man-yen around a mean of 349.9 (1 man-yen $=$ 10{,}000 JPY).

\begin{table}[h]
\centering
\begin{tabular}{lcc}
\toprule
Year & Mean separation rate & SD (10 seeds) \\
\midrule
2021 & 0.0850 & 0.0012 \\
2022 & 0.0800 & 0.0016 \\
2023 & 0.0765 & 0.0016 \\
2024 & 0.0755 & 0.0015 \\
\bottomrule
\end{tabular}
\caption{Dispersion of the simulated annual separation rate across ten
simulation seeds (backtest configuration of the main paper).}
\label{tab:sim_dispersion}
\end{table}

\paragraph{Bootstrap confidence intervals for the prediction benchmark}
We quantify sampling uncertainty in the separation benchmark (test wave 2024,
$n=27{,}638$) with a paired bootstrap: 1{,}000 resamples of the test set,
evaluating all six models on identical resamples. Table~\ref{tab:auc_ci}
reports 95\% percentile intervals for each model's ROC-AUC and for the AUC
advantage of XGBoost over every alternative. All gap intervals exclude zero:
XGBoost's advantage is statistically significant not only over the neural
sequence models and logistic regression but also over HistGradientBoosting.
The GRU and Transformer intervals overlap the current-state MLP almost
entirely, confirming that access to employment history does not improve
prediction on these short ($\le 8$-wave) sequences.

\begin{table}[h]
\centering
\begin{tabular}{lcc}
\toprule
Model & ROC-AUC [95\% CI] & Gap to XGBoost [95\% CI] \\
\midrule
XGBoost      & 0.728 [0.718, 0.738] & --- \\
HistGBM      & 0.709 [0.698, 0.720] & $+$0.019 [0.014, 0.024] \\
MLP          & 0.707 [0.695, 0.718] & $+$0.022 [0.015, 0.029] \\
GRU          & 0.705 [0.695, 0.716] & $+$0.023 [0.016, 0.029] \\
Transformer  & 0.704 [0.693, 0.714] & $+$0.025 [0.018, 0.031] \\
Logistic     & 0.657 [0.645, 0.668] & $+$0.072 [0.062, 0.082] \\
\bottomrule
\end{tabular}
\caption{Paired bootstrap (1{,}000 resamples) over the 2024 test set:
ROC-AUC with 95\% percentile intervals, and the AUC gap of XGBoost over each
alternative. Every gap interval excludes zero.}
\label{tab:auc_ci}
\end{table}

\section{Reproducibility Details}

\paragraph{Compute}
All experiments run on a single consumer laptop (Apple M1 CPU, 16\,GB RAM,
macOS~15.3; no GPU) under Python~3.9 with scikit-learn~1.6.1, XGBoost~2.1.4,
statsmodels~0.14.6, PyTorch~2.8.0 (CPU), SHAP~0.49.1, pandas~2.3.3 and
NumPy~2.0.2. The full pipeline --- panel construction through simulation and
the uncertainty analysis of Section~S1 --- completes in under one hour.

\paragraph{Runs and seeds}
Each result reported in the main paper is a single run with a fixed seed
(scikit-learn/XGBoost \texttt{random\_state}~$=$~42; PyTorch seed 0;
simulation seed 1). The dispersion analysis of Section~S1 uses simulation
seeds 1--10 and bootstrap seed 0.

\paragraph{Hyperparameters}
We performed no hyperparameter search; all values are fixed defaults chosen a
priori. \emph{Benchmark models} --- logistic regression: median imputation,
standardization, one-hot encoding, $C{=}1$, max 2{,}000 iterations, balanced
class weights; HistGradientBoosting: learning rate 0.05, 400 iterations,
$L_2{=}1.0$, native categorical splits, balanced class weights; XGBoost: 400
trees, learning rate 0.05, max depth 5, subsample 0.8, column subsample 0.8.
\emph{Transition sub-models}: HistGradientBoosting with learning rate 0.05,
300 iterations, $L_2{=}1.0$; \emph{no} class re-weighting in the classifiers
(calibration is required for the generative use, as discussed in the main
paper) and absolute-error loss in the wage regressors. \emph{Sequence models}
--- MLP: two hidden layers of width 64, dropout 0.2; GRU: hidden size 64;
Transformer: $d_{\mathrm{model}}{=}64$, 4 heads, 2 layers, feed-forward width
128, dropout 0.1; all trained with Adam (learning rate $10^{-3}$), batch size
512, 8 epochs, positively re-weighted binary cross-entropy, history length
$\le 8$. \emph{Microsimulation}: horizon 4 years; wage noise standard
deviation 30 man-yen.

\paragraph{Code and data}
The accompanying anonymized code archive contains the full pipeline
(harmonization map, loaders, feature builders, models, simulation, and the
scripts that generate every table and figure, including
\texttt{compute\_uncertainty.py} for Section~S1), a synthetic panel that runs
the entire pipeline end-to-end without restricted data, and the public
macroeconomic series (CPI; economy-wide and industry-level wage growth) with
source documentation. The JPSED microdata cannot be redistributed; qualified
researchers can obtain it from the Social Science Japan Data Archive, and the
released code reproduces all results once the microdata is in place.

\end{document}
